%%%%%%%%%%%%%%%%%%%%%%%%%%%%%%%%%%%%%%%%%
% Medium Length Professional CV
% LaTeX Template
% Version 2.0 (8/5/13)
%
% This template has been downloaded from:
% http://www.LaTeXTemplates.com
%
% Original author:
% Rishi Shah 
%
% Important note:
% This template requires the resume.cls file to be in the same directory as the
% .tex file. The resume.cls file provides the resume style used for structuring the
% document.
%
%%%%%%%%%%%%%%%%%%%%%%%%%%%%%%%%%%%%%%%%%

%----------------------------------------------------------------------------------------
%	PACKAGES AND OTHER DOCUMENT CONFIGURATIONS
%----------------------------------------------------------------------------------------

\documentclass{resume} % Use the custom resume.cls style
\usepackage[usenames, dvipsnames]{xcolor} % Optional: for more color names
\usepackage[colorlinks=true, urlcolor=blue, linkcolor=green, citecolor=red]{hyperref}

\usepackage{xurl}
\usepackage[left=0.75in,top=0.6in,right=0.75in,bottom=0.6in]{geometry} % Document margins
\newcommand{\tab}[1]{\hspace{.2667\textwidth}\rlap{#1}}
\newcommand{\itab}[1]{\hspace{0em}\rlap{#1}}
\name{Mihir Vinay Kulkarni} % Your name
\address{Researcher at the \href{https:www.ntnu.no}{Norwegian University of Science and Technology (NTNU)}} % Your address
%\address{123 Pleasant Lane \\ City, State 12345} % Your secondary addess (optional)

% \href{mailto:mihir.kulkarni@ntnu.no}{mihir.kulkarni@ntnu.no} $|$ 

\address{\href{mailto:mihirk284@gmail.com}{mihirk284@gmail.com} $|$ +(47)93933035 $|$ \href{https://www.github.com/mihirk284}{\bf{GitHub}} $|$
\href{https://scholar.google.com/citations?user=ioQvwcEAAAAJ&hl=en}{\bf{Google Scholar}}} 

%\href{https://sites.google.com/goa.bits-pilani.ac.in/mihir-kulkarni/}{\bf{Website}}} % Your phone number and email

\begin{document}

%----------------------------------------------------------------------------------------
%	EDUCATION SECTION
%----------------------------------------------------------------------------------------

\begin{rSection}{Education}
{\bf Ph.D. Engineering Cybernetics } \hfill {\em January 2022 - August 2025}
\\{\bf Norwegian University of Science and Technology}
\\ Department of Engineering Cybernetics\\
{\bf M.S. Computer Science and Engineering} \hfill {\em August 2020 - December 2021}
\\{\bf University of Nevada, Reno}
\\ Department of Computer Science and Engineering \hfill
(GPA: 3.875/4.0) \\
{\bf B.E. Mechanical Engineering} \hfill {\em August 2016 - July 2020} 
\\{\bf Birla Institute of Technology and Science, Pilani (Goa Campus)} 
\\ Department of Mechanical Engineering \hfill
(GPA: 8.66/10)

% % {\bf All India Senior School Certificate Examination} \hfill {\em July 2014 - June 2016} 
% % \\{\bf Delhi Public School, Navi Mumbai}
% \\ (Overall Percentage: 95\%)

\end{rSection}



%----------------------------------------------------------------------------------------
%	Publications
%----------------------------------------------------------------------------------------

\begin{rSection}{Journal Publications}
\begin{itemize}
    \item G. Malczyk, \textbf{M. Kulkarni} and K. Alexis, \textbf{``Semantically-Driven Deep Reinforcement Learning for Inspection Path Planning,''}, IEEE Robotics and Automation Letters. Available at \url{https://doi.org/10.1109/LRA.2025.3575331}.
    \item \textbf{M. Kulkarni}, W. Rehberg and K. Alexis, \textbf{``Aerial Gym Simulator: A Framework for Highly Parallelized Simulation of Aerial Robots''}, IEEE Robotics and Automation Letters. Available at \url{https://doi.org/10.1109/LRA.2025.3548507}.
    \item M. Tranzatto, M. Dharmadhikari, L. Bernreiter, M. Camurri, S. Khattak, F. Mascarich, P. Pfreundschuh, D. Wisth, S. Zimmermann, \textbf{M. Kulkarni}, V. Reijgwart, B. Casseau, T. Homberger, P. De Petris, L. Ott, W. Tubby, G. Waibel, H. Nguyen, C. Cadena, R. Buchanan, L. Wellhausen, N. Khedekar, O. Andersson, L. Zhang, T. Miki, T. Dang, M. Mattamala, M. Montenegro, K. Meyer, X. Wu, A. Briod, M. Mueller, M. Fallon, R. Siegwart, M. Hutter, K. Alexis, \textbf{``Team CERBERUS Wins the DARPA Subterranean Challenge: Technical Overview and Lessons Learned''}, Field Robotics. Available at \url{https://doi.org/10.48550/arXiv.2207.04914}.
    \item M. Tranzatto, T. Miki, M. Dharmadhikari, L. Bernreiter, \textbf{M. Kulkarni}, F. Mascarich, O. Andersson, S. Khattak, M. Hutter, R. Siegwart, K. Alexis, \textbf{``CERBERUS in the DARPA Subterranean Challenge''} Science Robotics. Available at \url{https://www.science.org/doi/abs/10.1126/scirobotics.abp9742}.
    \item F. Mascarich, \textbf{M. Kulkarni}, P. de Petris, T. Wilson, K. Alexis, \textbf{``Autonomous Mapping and Spectroscopic Analysis of Distributed Radiation Fields using Aerial Robots"}, Autonomous Robots. Available at \url{https://doi.org/10.1007/s10514-022-10064-7}.
    \item M. Tranzatto, F. Mascarich, L. Bernreiter, C. Godinho, M. Camurri, S. Khattak, T. Dang, V. Reijgwart, J. Loje, D. Wisth, S. Zimmermann, H. Nguyen, M. Fehr, L. Solanka, R. Buchanan, M. Bjelonic, N. Khedekar, M. Valceschini, F. Jenelten, M. Dharmadhikari, T. Homberger, P. De Petris, L. Wellhausen, \textbf{M. Kulkarni}, T. Miki, S. Hirsch, M. Montenegro, C. Papachristos, F. Tresoldi, J. Carius, G. Valsecchi, J. Lee, K. Meyer, X. Wu, J. Nieto, A. Smith, M. Hutter, R. Siegwart, M. Mueller, Ma. Fallon, K. Alexis, \textbf{``CERBERUS: Autonomous Legged and Aerial Robotic Exploration in the Tunnel and Urban Circuits of the DARPA Subterranean Challenge''}, Field Robotics. Available at \url{https://doi.org/10.55417/fr.2022011}.\\ 
\end{itemize}
\end{rSection}

\newpage
\begin{rSection}{Conference Publications}
\begin{itemize}
    \item M. Kulkarni, M. Dharmadhikari, N. Khedekar, M. Nissov, M. Singh, P. Weiss andK. Alexis, 2025. \textbf{``UniPilot: Enabling GPS-Denied Autonomy Across Embodiments''}. IEEE International Conference on Advanced Robotics (ICAR) 2025. Available at: \url{https://doi.org/10.48550/arXiv.2509.11793}.
    \item M. Harms, \textbf{M. Kulkarni}, N. Khedekar, M. Jacquet, K. Alexis. \textbf{``Neural Control Barrier Functions for Safe Navigation''}. IEEE/RSJ International Conference on Intelligent Robots and Systems (IROS) 2024. Available at: \url{https://doi.org/10.1109/IROS58592.2024.10802694}.
    \item \textbf{M. Kulkarni} and K. Alexis. \textbf{``Reinforcement Learning for Collision-free Flight Exploiting Deep Collision Encoding''}. IEEE International Conference on Robotics and Automation (ICRA) 2024. Available at \url{https://doi.org/10.48550/arXiv.2402.03947}.
    \item M. Dharmadhikari, P. De Petris, \textbf{M. Kulkarni}, N. Khedekar, H. Nguyen, A.E. Stene, E. Sjøvold, K. Solheim, Bente Gussiaas, and Kostas Alexis. \textbf{``Autonomous Exploration and General Visual Inspection of Ship Ballast Water Tanks using Aerial Robots.''}, IEEE International Conference on Advanced Robotics (ICAR) 2023. \textit{\textbf{Winner - Best Paper Award}}. Available at \url{https://doi.org/10.48550/arXiv.2311.03838}.
    \item \textbf{M. Kulkarni} and K. Alexis, \textbf{``Task-driven Compression for Collision Encoding based on Depth Images''}. International Symposium on Visual Computing (ISVC) 2023. Available at \url{https://doi.org/10.48550/arXiv.2309.05289}.
    \item \textbf{M. Kulkarni}, H. Nguyen, and K. Alexis. \textbf{``Semantically-enhanced Deep Collision Prediction for Autonomous Navigation using Aerial Robots''}. IEEE/RSJ International Conference on Intelligent Robots and Systems (IROS) 2023. Available at \url{https://doi.org/10.48550/arXiv.2307.11522}.
    \item M. Dharmadhikari, P. De Petris, H. Nguyen, \textbf{M. Kulkarni}, N. Khedekar and K. Alexis, \textbf{``Manhole Detection and Traversal for Exploration of Ballast Water Tanks using Micro Aerial Vehicles''}, International Conference on Unmanned Aircraft Systems (ICUAS) 2023, Available at \url{https://doi.org/10.1109/ICUAS57906.2023.10156214}.
    \item N. Khedekar, \textbf{M. Kulkarni} and K. Alexis, \textbf{``MIMOSA: A Multi-Modal SLAM Framework for Resilient Autonomy against Sensor Degradation''}, IEEE/RSJ International Conference on Intelligent Robots and Systems (IROS) 2022. Available at \url{https://doi.org/10.1109/IROS47612.2022.9981108}.
    \item P. De Petris, H. Nguyen, M. Dharmadhikari, \textbf{M. Kulkarni}, N. Khedekar, F. Mascarich, and K. Alexis. \textbf{``RMF-Owl: A Collision-Tolerant Flying Robot for Autonomous Subterranean Exploration.''}, International Conference on Unmanned Aircraft Systems (ICUAS) 2022. Available at \url{https://doi.org/10.1109/ICUAS54217.2022.9836115}.
    \item \textbf{M. Kulkarni}, M. Dharmadhikari, M. Tranzatto, S. Zimmermann, V. Reijgwart, P. De Petris, H. Nguyen, N. Khedekar, C. Papachristos, L. Ott, R. Siegwart, M. Hutter, and K. Alexis, \textbf{``Autonomous Teamed Exploration of Subterranean Environments using Legged and Aerial Robots''}, IEEE International Conference on Robotics and Automation (ICRA) 2022. Available at \url{https://doi.org/10.1109/ICRA46639.2022.9812401}. \textit{\textbf{Winner - Outstanding Deployed Systems Paper Award}}.
    \item P. De Petris, H. Nguyen, \textbf{M. Kulkarni}, F. Mascarich and K. Alexis, \textbf{``Resilient Collision-tolerant Navigation in Confined Environments''}, 2021 IEEE International Conference on Robotics and Automation (ICRA), 2021. Available at \url{https://doi.org/10.1109/ICRA48506.2021.9561999}.
    \item \textbf{M. Kulkarni}, H. Nguyen, and K. Alexis, \textbf{``The Reconfigurable Aerial Robotic Chain: Shape and Motion Planning''}. IFAC World Congress, 2020. Available at \url{https://doi.org/10.1016/j.ifacol.2020.12.2383}.
\end{itemize}
\end{rSection}

\begin{rSection}{Whitepapers}
\begin{itemize}
    \item M. Mittal et. al. \textbf{``Isaac Lab: A GPU Accelerated Simulation Framework For Multi-Modal Robot Learning''}. Available at: \url{https://research.nvidia.com/publication/2025-09_isaac-lab-gpu-accelerated-simulation-framework-multi-modal-robot-learning}.
    \item M. Dharmadhikari et. al. \textbf{``The Unified Autonomy Stack: Toward a Blueprint for Generalizable Robot Autonomy''}. Available at: \url{https://ntnu-arl.github.io/unified_autonomy_stack/report/UnifiedAutonomyStack-Outline/}
\end{itemize}
\end{rSection}


\begin{rSection}{Book Chapters}
\begin{itemize}
    \item \textbf{Mihir Kulkarni}, Brady Moon, Sebastian Scherer, Kostas Alexis, \textbf{``Aerial Field Robotics''}, Encyclopedia of Robotics. Available at \url{https://doi.org/10.1007/978-3-642-41610-1_221-1}.

\end{itemize}
\end{rSection}


\begin{rSection}{Talks and Lectures}
\begin{itemize}
    \item PX4 Developer Summit 2025 - \textbf{From Pixels To Propellers: Sim2Real Control and Vision-based Navigation}. \href{https://www.youtube.com/watch?v=r1ws_NJmcJk}{Link}
    \item Invited Lecture - WPI RBE-595-F02-ST: \textbf{Reinforcement learning for control and navigation of aerial robots}
    \item Learning-oriented Simulation for Aerial Robots, SSCI 2025 - \textbf{Tutorial: Aerial Gym 2.0: Isaac Gym-based Massively Parallelized Simulation for Efficient Aerial Robot Learning}
\end{itemize}
    
\end{rSection}


\begin{rSection}{Awards and Achievements}
\begin{itemize}
    \item Best Paper Award \hfill  IEEE ICAR 2023
    \item Outstanding Deployed Systems Paper Award \hfill IEEE ICRA 2022
    \item Certificate of Special Recognition \hfill United States Senate (2021)
    \item Winner - Prize Round \hfill DARPA Subterranean Challenge (2021)
\end{itemize}
\end{rSection}


\begin{rSection}{Open Source Contributions}
\begin{itemize}
    \item Unified Autonomy Stack - a field-tested autonomy architecture commanding a diverse set of robots. \href{https://github.com/ntnu-arl/unified_autonomy_stack}{\textbf{GitHub}} \href{https://ntnu-arl.github.io/unified_autonomy_stack}{\textbf{Website}}
    \item Aerial Gym Simulator - massively parallelized aerial robot simulator based on NVIDIA Isaac Gym. \href{https://github.com/ntnu-arl/aerial_gym_simulator}{\textbf{GitHub}} \href{https://ntnu-arl.github.io/aerial_gym_simulator}{\textbf{Website}}.
    \item Semantically-enhanced Variational Autoencoder. \href{https://github.com/ntnu-arl/sevae}{\textbf{GitHub}}.
    \item GSOC 2020: Sensor Data Visualization - Open Robotics. \href{https://community.gazebosim.org/t/gsoc-2020-sensor-data-visualization/638}{\textbf{Website}}.
    \item Simulation Models - Team CERBERUS - DARPA Subterranean Challenge Simulator. \href{https://github.com/osrf/subt/tree/final_event/submitted_models}{\textbf{GitHub}}.
    \item SuperMegaBot Simulator - Team CERBERUS Roving Robot \href{https://github.com/ntnu-arl/smb_simulator}{\textbf{GitHub}}.
\end{itemize} 
\end{rSection}


%----------------------------------------------------------------------------------------
%	Key Skills
%----------------------------------------------------------------------------------------


\begin{rSection}{Programs, Internships and Experience}
\small{{{\bf Nordic Probabilistic AI School.} \hfill \small{{\em June 2023}}}}\\
\small{{{\bf Google Summer of Code 2020, OpenRobotics.} \hfill \small{{\em May 2020 - September 2020}}}}\\
\small{{\bf Visiting Scholar, Autonomous Robots Lab, University of Nevada, Reno.} \hfill \small{{\em June 2019 - January 2020}}}\\
{\bf Summer Research Intern, CSIR-CEERI Pilani, India.} \hfill \small{{\em May 2018 - July 2018}}
\end{rSection}

\begin{rSection}{Technical Skills and Proficiencies}
{\bf{Programming Languages}} - C, C++, Python\\
{\bf{Mechanical Design}} - SOLIDWORKS, PTC Creo, Autodesk Fusion 360\\
{\bf{Proficient in ROS, ROS 2, Gazebo Classic and (Ignition) Gazebo}}\\
{\bf{Experienced in mechanical design, hardware and sensor integration and full-stack autonomy integration for high-performance mobile robotic platforms.}}
\end{rSection}

\begin{rSection} {Print and Digital Media}
\begin{itemize}
    \item Multiple features in IEEE Spectrum - Video Friday
    \begin{itemize}
        \item Reinforcement Learning for Collision-free Flight Exploiting Deep Collision Encoding. \url{https://www.youtube.com/watch?v=gPrT21sbpTY}
        \item{Autonomous Teamed Exploration of Subterranean Environments using Legged and Aerial Robots. \url{https://spectrum.ieee.org/video-friday-perseverance-autonomy}} 
        \item Semantically-enhanced Deep Collision Prediction for Autonomous Navigation using Aerial Robots. \url{https://spectrum.ieee.org/video-friday-resilient-bugbots}.
    \end{itemize}
    
    \item Featured in The Washington Post Magazine ``The Pentagon’s \$82 Million Super Bowl of Robots'' \url{https://www.washingtonpost.com/magazine/2021/11/10/darpa-robot-competition/}.
    \item Featured in the Teknisk Ukeblad article ``Seier for NTNU-basert robotmiljø: Bedre enn både Nasa og MIT'' \url{https://www.tu.no/artikler/seier-for-ntnu-basert-robotmiljo-bedre-enn-bade-nasa-og-mit/513808}.
    \item GazeboSim Community: GSOC 2020: Sensor Data Visualisation \url{https://community.gazebosim.org/t/gsoc-2020-sensor-data-visualization/638}.
    \item BITS R\&D Post: Thesis at The University of Nevada, Reno \url{https://bitsrnd.wordpress.com/2020/04/21/mihir-kulkarnis-thesis-at-university-of-nevada-reno/}.
\end{itemize}
\end{rSection}


%	ACADEMIC COURSES
%----------------------------------------------------------------------------------------

% \begin{rSection}{Academic Courses} 
%  Machine Design and Drawing, Kinematics and Dynamics of Machines, Nonlinear Dynamics and Chaos Theory, Computer Aided Design, Mechanical Vibrations, Advanced Mechanics of Solids, Fluid Mechanics, Control Systems, Computer Programming, Software Development for Portable Devices, Analysis of Algorithms, Deep Learning, Convex Optimization, Autonomous Mobile Manipulation.
% \end{rSection}



%----------------------------------------------------------------------------------------
%	Professional Experience
%----------------------------------------------------------------------------------------

% \begin{rSection}{Professional Experience}
% {\bf Visiting Scholar} \hfill \small{{\em June 2019 - December 2019}}\\
% Autonomous Robots Lab, University of Nevada, Reno


% {\bf Academic Research Assistant} \hfill \small{{\em January 2019 - May 2019}}\\
% Department of Electrical and Electronics Engineering, BITS Pilani

% {\bf Academic Research Assistant} \hfill \small{{\em August 2018 - December 2018}}\\
% Department of Mechanical Engineering, BITS Pilani

% {\bf Summer Research Intern, Cyber-Physical Systems Division} \hfill \small{{\em May 2018 - July 2018}}\\
% Central Electronics Engineering Research Institute, Pilani
% \end{rSection}








% %--------------------------------------------------------------------------------
% %    Research Experience
% %-----------------------------------------------------------------------------------------------

% \begin{rSection}{Research Experience}

% {\bf COHORT: CoOperation for HeterOgeneous Robot Teams}
% \\\textcolor{darkgray}{\small{\bf Advisor: Prof. Dr. Kostas Alexis}}
% \\ \textcolor{darkgray}{ August 2020 - December 2020 $|$ Autonomous Robots Lab, University of Nevada, Reno}
% \\ Developed a framework for coordination of heterogeneous teams of autonomous robots. A central base-station collects parts of a volumetric map from individual robots and build a global occupancy map. Using neighbor search on the discretized environment, a neighbor search allows for finding local forntiers in the parts of volumetric maps, that are combined globally to identify frontiers for the environment. The framework enables repositioning of the robots to global frontiers leading to more efficient exploration.

% {\bf Trajectory and Shape Planning of a Reconfigurable Aerial Robotic System-of-Systems}
% \\\textcolor{darkgray}{\small{\bf Bachelor's Thesis}}
% \\\textcolor{darkgray}{\small{\bf Advisor: Prof. Dr. Kostas Alexis}}
% \\ \textcolor{darkgray}{ June 2019 to December 2019 $|$ Autonomous Robots Lab, University of Nevada, Reno}
% \\ Developed a motion planning algorithm and probabilistic roadmap-based methods for the translation and shape of the reconfigurable system. The de-coupled planning approach uses lesser sampled points compared to traditional full-state PRM methods, which is an order of magnitude faster and works well for systems with high degrees of freedom.

% % {\bf Autonomous Drone System}
% % \\ \textcolor{darkgray}{\small{\bf Dr. Shibu Clement, Mechanical Engineering}}
% % \\ \textcolor{darkgray}{August 2018 to December 2018 $|$ BITS Pilani, Goa}
% % \\ Designed and interfaced an aerial robot platform to perform surveillance of a given area, while avoiding obstacles. Performed autonomous flight to collect visual data of a given area. Designed a browser-based ground-control station to control the quadrotor.\\

% {\bf Path Planning for Autonomous Quadrotor}
% \\\textcolor{darkgray}{\small{\bf Summer Internship}}
% \\ \textcolor{darkgray}{\small{\bf Advisor: Dr. Bhausaheb Botre}}
% \\ \textcolor{darkgray}{May 2018 to July 2018 $|$ Central Electronics Engineering Research Institute, Pilani}
% \\Indigenously developed aerial platform, fully equipped with onboard sensors, capable of navigating across obstacle-laden regions to inspect power lines. Also set up automated pipelines for flight motions, data collection, and visualization. Path planning was performed using the Potential Field based method with modifications to overcome local minima. The system was based on ROS and the simulations were performed in Gazebo. Conducted virtual and physical validation of the system.


% {\bf Modelling and Simulation of Chaotic System: Forced-Damped Duffing Oscillator}
% \\ \textcolor{darkgray}{\small{\bf Dr. Gaurav Dar, Department of Physics}}
% \\ \textcolor{darkgray}{August 2017 to December 2017 $|$ BITS Pilani, Goa}
% \\Designed and evaluated a Forced-Damped Oscillator to model chaotic behaviour in nonlinear stiffening of springs, bending of beams and electronic circuits. Performed evaluation studies in MATLAB and Simulink to model the behaviour and response of the oscillator.Also investigated behaviour of forced Van-Der-Pol Oscillator.

% \end{rSection}

%--------------------------------------------------------------------------------
%    Projects And Seminars
%-----------------------------------------------------------------------------------------------
% \begin{rSection}{Key Projects}

% {\bf Team CERBERUS - DARPA Subterranean Challenge - Winning Team}\\
% I was a part of Team CERBERUS, participating in the DARPA Subterranean Challenge. The challenge focuses on performing search and rescue missions using autonomous robots in underground settings, in GPS-denied, degraded environment and sensing conditions. I worked on our team's ground robot, equipped with multiple LIDAR sensors, thermal cameras, depth and colour cameras and multiple network interfaces. I worked on the motion planning and controls of the skid-steer drive robot and further designed a robust sensor head to accommodate all the above sensors on the robot. I have also worked on the artifact detection pipeline for the team that included the training and deploying the inference pipeline of the team's object detector using YOLO-v3. Team CERBERUS won the first position during the Final Event of the DARPA Subterreanean Challenge. Please see following links for more information - \href{http://www.subt-cerberus.org/}{\textbf{Team Website}}.



% {\bf Hyperloop India - SpaceX Hyperloop Pod Competition 2019}\\
% Team Lead for the Electronics Team in Hyperloop India. Worked on the electronic subsystem for the hyperloop pod, system design, selection and interfacing of electronics components. Coordinated with the mechanical design team for integration. Gained major learning during The SpaceX - Hyperloop Pod competition.


% As the Team Lead for the Electronics Team in Hyperloop India, I worked on the electronics subsystem for the pod. I worked on the system design and the selection and interfacing of electronics components. With my knowledge of mechanical design, I coordinated with the design team as well. The Hyperloop pod competition was a major part of my learning experience.\\


% {\bf Nano-Quadrotor on a Printed Circuit Board frame}\\
% Designed a PCB based Nano Quadrotor frame and built a transmitter using ATmega328P based microcontroller boards. Designed communication system for the drone using nRF24L01+ modules. Tested the hardware using PID attitude controller. Currently working to develop the platform for using Model Predictive Controller and enabling compatibility with a base-station using ROS. Please refer to \href{https://sites.google.com/goa.bits-pilani.ac.in/mihir-kulkarni/#h.p_Sj1AjuVjXMhR}{\bf{this link}} for additional information and images of the prototype.



% {\bf Simulation of Multi-Agent Quadrotor Swarms}\\
% Developed the pipeline for controlling multiple quadrotors for a flying light display. Worked on task allocation for individual agents in a swarm. Simulated a multi-agent system in Gazebo to dynamically include new drones without affecting existing formations The robots are controlled using MAVROS. Both velocity based control and position based control is implemented on each of the robots. The videos of the simulations can be found \href{https://sites.google.com/goa.bits-pilani.ac.in/mihir-kulkarni/#h.p_Ey2yJ-dnbUtq}{\bf{here}}.


% Autonomous aerial platforms coordinating with other identical agents can lead to completion of complex tasks effectively. I worked on the task allocation for the individual agents in a swarm. I have simulated a multi agent system in Gazebo to dynamically include new drones at dynamically without affecting existing formations.\\


% I worked as the Electronics Team Lead for the competition. I worked on the design of the electrical and electronics systems of the pod and worked with the software team to optimise the design as per the competition specifications. I gained a lot of knowledge and experience working on the hardware, and I was comfortable interfacing the various components involved in the Hyperloop Pod. With my knowledge of Mechanical engineering, I was able to coordinate with the mechanical team to integrate the two subsystems seamlessly.\\



% {\bf Design and Implementation of a test bed for swarm algorithms}\\
% The project aim to develop an autonomous terrestrial robot capable of navigation in human environments, capable of planning around dynamic obstacles to deliver parcels to offices and rooms. I have worked on the hardware interfacing and of the robot and am working on AR tag detection and communication and on the drive system using the omni-wheel based holonomic control. I am also working on the navigation team for dynamic path planning of the robot. The team is currently exploring ways to navigate using combinations of different kinds of non-visual sensors.\\



% {\bf DRDO Robotics and Unmanned Systems Exposition (DRUSE)}\\
% Team member to design a high-performance multirotor system for defense purposes. The drone is designed to have a high payload capacity and an extended range. Worked on the mechanical design of camera gimbals and implemented optimal trajectory-generation and tracking.


% I worked as a part of a team to design a high performance multirotor system for defense purposes. The drone is designed to have a high payload capacity and extended range. I worked on the mechanical design of the camera gimbals and implemented optimal trajectory-generation and tracking.\\


% {\bf Autonomous Delivery Bot}\\
% The project aims to develop a robot capable of navigation in dynamic environments. I have worked on the hardware interfacing and of the robot and developed the motion planning solution using a 2D lidar and an RRT* planning algorithm for holonomic drive.\\

% {\bf Wearable Gesture Based HCI Controller}\\
% Developed and fabricated a wireless, wearable gesture controller that detects pose and movement of the human arm. Applied machine learning algorithms to identify and classify motions. It can be used to control a variety of devices ranging from remote controlled cars, UAVs, to home appliances. Worked on the interfacing of sensors, data collection, visualization and real-time control.




% I have developed and fabricated a wireless, wearable gesture controller that detects pose and the movement of the human arm and
% applies machine learning algorithms to classify and identify the motions. It can be used to control a variety of devices ranging from remote controlled cars, drones, to home appliances. I have worked on the sensor interfacing, data collection, visualization, and real-time control.\\

% {\bf Non-Visual Crowd Counting using RSSI}\\
% Non-visual counting of people in an environment is an efficient, sensor based method to detect the number of people passing through two sensor modules. The RSSI of the WiFi signal  changes depending on the number, size and speed of obstacles passing between the line of sight of two sensor nodes. The changes in the RSSI signal are detected and the number of people passing can be computed in quick manner. I set up data-collection and filering pipelines and was responsible for interfacing the sensors to collect and process data in real time.\\

% \end{rSection}


%----------------------------------------------------------------------------------------
%	Conferences Refereed
%----------------------------------------------------------------------------------------

\begin{rSection}{Journals and Conferences Refereed}

{\bf Journals:} IEEE RA-L, IEEE RAM, IJRR, IEEE T-RO, IEEE T-FR \\
{\bf Conferences:} IEEE ICRA, IEEE/RSJ IROS, IEEE ICUAS, IEEE ICAR
\end{rSection}





%----------------------------------------------------------------------------------------
%	POR
%----------------------------------------------------------------------------------------





\begin{rSection}{Positions of Responsibility}
\small{{{\bf Chief Coordinator, Aerodynamics Club, BITS Pilani} \hfill \small{{\em March 2018 - May 2019}}}}\\
\small{{\bf Sub-Coordinator, Electronics and Robotics Club, BITS Pilani} \hfill \small{{\em April 2018 - May 2019}}}\\
{\bf Electronics Team Lead, Hyperloop India} \hfill \small{{\em March 2018 - January 2019}}


\end{rSection}



%----------------------------------------------------------------------------------------
%	Teaching Experience
%----------------------------------------------------------------------------------------

\begin{rSection}{Teaching and Research Experience}
\small{\bf Graduate Research Assistant} - Department of CSE, UNR \hfill \small{{\em Jan 2021 - Dec 2021}} \\
\small{\bf Teaching Assistant} - Computer Programming (Department of CSE, BITS Pilani) \hfill \small{{\em Jan 2020 - May 2020}} \\
\small{\bf Instructor} - Intermediate Robotics (Center for Technical Education, BITS Pilani) \hfill \small{{\em Jan 2019 - May 2019}} \\
{\bf Instructor} - Introduction to Robotics (Center for Technical Education, BITS Pilani) \hfill \small{{\em Aug 2018 - Dec 2018}}\\
\end{rSection}

% \begin{rSection}{Referees}
% \small{\bf Prof. Dr. Kostas Alexis} - Dept. of Engineering Cybernetics, NTNU \hfill \small{konstantinos.alexis@ntnu.no} \\
% \small{\bf Prof. Dr. Christos Papachristos} - Dept. of Computer Science, UNR \hfill \small{cpapachristos@unr.edu} \\
% \small{\bf Prof. Dr. Shibu Clement} - Dept. of Mechanical Engineering, BITS Pilani \hfill \small{shibu@goa.bits-pilani.ac.in} \\



\end{document}
